\vspace{-3.6em}
\begin{abstract}
\footnotesize
\linespread{0.92}\selectfont
\noindent
\textbf{Objetivo.} Identificar quem está exposto à inteligência artificial (IA) no mercado de trabalho brasileiro, marcado por ampla informalidade contratual, e por quais mecanismos econômicos essa exposição se distribui. \textbf{Métodos.} Desenvolvo um modelo de designação de tarefas em que a IA individual está universalmente disponível, mas a IA corporativa exige capital organizacional restrito ao setor formal. Testo as previsões do modelo nos microdados individuais da Pesquisa Nacional por Amostra de Domicílios Contínua (PNAD Contínua/IBGE, 2026Q1), combinados ao índice de exposição ocupacional à IA de \citet{eloundou2024gpts}, em escala de 0 a 10. O gradiente escolaridade--exposição e a mediação educacional da associação exposição--informalidade são estimados por regressões ponderadas no nível individual ($N = 227.629$; 122 ocupações), com controles mincerianos completos, erros-padrão agrupados por ocupação e robustez à seleção em não observáveis avaliada pelo limiar $\delta$ de Oster. \textbf{Resultados.} 35\% da força de trabalho ocupada (35{,}7 milhões de pessoas) atua em ocupações de exposição quase nula ($\theta \le 2$). Cada ano adicional de escolaridade associa-se a 0{,}23 ponto de exposição (erro-padrão 0{,}03; $p < 0{,}001$), gradiente que resiste à seleção em não observáveis ($\delta = 1{,}93 > 1$). A associação negativa entre exposição e informalidade ($-6{,}23$ pontos percentuais) atenua-se em 37\% após o controle educacional ($-3{,}91$ pontos percentuais; $p < 0{,}001$), evidenciando mediação parcial consistente com o canal do capital organizacional. \textbf{Conclusões.} A exposição à IA concentra-se no núcleo formal e escolarizado, e a informalidade contratual opera como amortecedor estrutural da automação corporativa por limitar a extração de ganhos organizacionais de produtividade --- e não o acesso às ferramentas. Políticas de capacitação e transição tecnológica devem priorizar o estrato médio formalizado.\par
\vspace{0.12em}
\noindent \textbf{Palavras-chave:} Inteligência Artificial, Mercado de Trabalho, Informalidade, Capital Organizacional, Pareamento Ocupacional, PNAD Contínua. \textbf{Classificação JEL:} J23, J24, O33, O17, J46.\par
\vspace{0.12em}
\noindent \rule{\linewidth}{0.4pt}\par
\vspace{0.12em}
\noindent \textbf{Abstract:}
\textbf{Objective.} To identify who is exposed to artificial intelligence (AI) in the Brazilian labor market, characterized by widespread contractual informality, and through which economic mechanisms that exposure is distributed. \textbf{Methods.} I develop a task-assignment model in which individual AI is universally available, but corporate AI requires organizational capital restricted to the formal sector. I test the model's predictions on individual-level microdata from the Brazilian Continuous National Household Sample Survey (PNAD Contínua/IBGE, 2026Q1), linked to the occupational AI exposure index of \citet{eloundou2024gpts}, on a 0--10 scale. The schooling--exposure gradient and the educational mediation of the exposure--informality association are estimated with weighted individual-level regressions ($N = 227{,}629$; 122 occupations), full Mincerian controls, occupation-clustered standard errors, and robustness to selection on unobservables assessed by Oster's $\delta$ threshold. \textbf{Results.} 35\% of the employed workforce (35.7 million workers) holds near-zero-exposure occupations ($\theta \le 2$). Each additional year of schooling is associated with a 0.23-point increase in AI exposure (SE 0.03; $p < 0.001$), a gradient robust to selection on unobservables ($\delta = 1.93 > 1$). The negative exposure--informality association ($-6.23$ percentage points) attenuates by 37\% after controlling for education ($-3.91$ percentage points; $p < 0.001$), indicating partial mediation consistent with the organizational-capital channel. \textbf{Conclusions.} AI exposure concentrates in the formal, educated core, and contractual informality operates as a structural buffer against corporate automation by limiting the extraction of organizational productivity gains --- not access to the tools. Upskilling and technological-transition programs should prioritize the mid-skilled formal stratum.\par
\vspace{0.12em}
\noindent \textbf{Keywords:} Artificial Intelligence, Labor Markets, Informality, Organizational Capital, Occupational Sorting, PNAD Contínua. \textbf{JEL Codes:} J23, J24, O33, O17, J46.\par
\end{abstract}
\newpage

